\section{Заключение}
\label{sec:Chapter5}

В~настоящей работе исследованы методы автоматической классификации
нестационарных биосигналов~--- электрокардиограмм~(ЭКГ)~---
с~использованием алгоритмов машинного обучения.
На~основе базы~данных MIT-BIH Arrhythmia Database
($N = 100\,025$ кардиоциклов пяти классов аритмий)
реализован полный конвейер обработки данных, обучены и~сравнены
три модели, а~также проведен анализ их~устойчивости к~шуму.

\subsection*{Основные результаты}

\begin{enumerate}

\item \textbf{Разработан конвейер предобработки данных.}
Реализовано извлечение кардиоциклов в~виде окон
длиной~$240$~отсчетов ($100$ до~и $140$ после R-пика),
стратифицированное разбиение 80/20, стандартная нормализация
для~нейронных сетей.
Для~борьбы с~дисбалансом классов (соотношение норма~:~предсердная
экстрасистола~$\approx 30\!:\!1$) применена взвешенная функция
потерь (формула~\eqref{eq:wce}) и~стратегия
\texttt{balanced\_subsample} для~Random Forest.

\item \textbf{Обучены три модели классификации.}
На~чистых данных достигнуты следующие результаты
(таблица~\ref{tab:clean_results}):

\begin{itemize}
    \item \textbf{1D~CNN}: Accuracy~$= 0.9917$,
          Macro~F1~$= 0.974$~--- наилучший результат.
          Recall класса~A (предсердная экстрасистола)~$= 0.92$.
    \item \textbf{ResNet1D}: Accuracy~$= 0.985$,
          Macro~F1~$= 0.953$.
    \item \textbf{Random Forest}: Accuracy~$= 0.985$,
          Macro~F1~$= 0.950$.
          Recall класса~A~$= 0.70$~--- существенно ниже
          нейронных сетей.
\end{itemize}

\item \textbf{Проведен вейвлет-анализ ЭКГ-сигналов.}
С~помощью дискретного вейвлет-преобразования (db4, 4~уровня)
визуализировано частотно-временно\'{е} строение кардиоциклов
каждого класса.
Проверена гипотеза о~добавлении вейвлет-признаков к~вектору
временны\'{х} отсчетов: расширение пространства признаков
с~$240$ до~$265$ не~улучшило качество Random Forest
(Macro~F1: $0.9479$ против~$0.9496$), что~объясняется
способностью ансамбля деревьев неявно улавливать
частотно-временны\'{е} паттерны из~сырого сигнала.

\item \textbf{Исследована устойчивость к~четырем типам помех.}
Выявлены характерные профили устойчивости:

\begin{itemize}
    \item CNN~--- наилучшая устойчивость к~дрейфу изолинии
    (Macro~F1~$= 0.974$) и~мышечному шуму ($0.960$):
    BatchNorm нейтрализует медленные тренды.
    \item Random Forest~--- наилучшая устойчивость к~гауссову
    ($0.872$) и~сетевому ($0.911$) шуму:
    пороговые разбиения нечувствительны к~равномерно
    распределенным и~периодическим аддитивным помехам.
    \item ResNet1D~--- наименее устойчив к~сетевой помехе~50~Гц
    (Macro~F1~$= 0.640$): глобальный усредняющий пулинг
    усиливает периодические артефакты в~признаковом
    пространстве.
\end{itemize}

\item \textbf{Достоверность подтверждена кросс-валидацией.}
Пятикратная стратифицированная кросс-валидация показала:
Random Forest~--- наиболее стабильная модель
(std Macro~F1~$= 0.002$); 1D~CNN~--- средняя стабильность
(std~$= 0.020$); ResNet1D~--- высокая дисперсия по~фолдам
(std~$= 0.053$) вследствие нехватки эпох обучения
при~CV-режиме.

\end{enumerate}

\subsection*{Выводы}

Все три~модели превысили порог Macro~F1~$\geq 0.90$
на~тестовой выборке и~порог recall(A)~$\geq 0.65$,
что~соответствует требованиям постановки задачи
(раздел~\ref{sec:Chapter1}).

Если приоритет~--- \textit{максимальная точность},
рекомендуется \textbf{1D~CNN}:
она лидирует по~всем метрикам на~чистых данных
и~устойчива к~наиболее распространенным помехам
(дрейф изолинии, мышечный шум).

Если приоритет~--- \textit{воспроизводимость и~интерпретируемость},
предпочтительнее \textbf{Random Forest}:
минимальная дисперсия по~фолдам, возможность анализа
важности признаков и~отсутствие необходимости
в~дорогостоящем обучении.

\subsection*{Направления дальнейших исследований}

\begin{itemize}
    \item \textbf{Вейвлет-шумоподавление как стадия предобработки.}
    Применение мягкой пороговой обработки коэффициентов ДВП
    (метод Донохо--Джонстона~\cite{donoho1994}) перед подачей
    сигнала на~классификатор может повысить устойчивость
    всех~трех моделей к~гауссову и~сетевому шуму.

    \item \textbf{Увеличение числа эпох ResNet1D в~кросс-валидации.}
    Полное обучение~(30~эпох) дает Macro~F1~$= 0.953$;
    при~CV использовалось лишь~$10$~эпох~---
    что~объясняет нестабильность. Рекомендуется тестировать
    с~$25$--$30$~эпохами.

    \item \textbf{Расширение набора классов.}
    MIT-BIH содержит более~$15$ типов аномалий;
    включение дополнительных классов (например, блокады
    1-й и~2-й степени, мерцательная аритмия) приблизит
    задачу к~реальной клинической практике.

    \item \textbf{Персонализированные модели.}
    Записи разных пациентов существенно различаются по~амплитуде
    и~форме ЭКГ; обучение с~адаптацией к~конкретному
    пациенту~(patient-specific fine-tuning) является
    перспективным направлением.

    \item \textbf{CWT-признаки для нейронных сетей.}
    Глубокие сверточные сети уже показали высокое качество
    при~прямой работе с~одномерным ЭКГ-сигналом~\cite{acharya2017};
    перспективное расширение~--- подача скаллограмм непрерывного
    вейвлет-преобразования~(CWT) в~двумерную~CNN как~изображений
    «частота~--- время».
\end{itemize}
